\documentclass[sigconf,balance=false]{acmart}
\usepackage{popets}

% Copyright
\setcopyright{popets}
\copyrightyear{YYYY}

% Issue info
\acmYear{YYYY}
\acmVolume{YYYY}
\acmNumber{X}
\acmDOI{XXXXXXX.XXXXXXX}
\acmISBN{}
\acmConference{Proceedings on Privacy Enhancing Technologies}
\settopmatter{printacmref=false,printccs=false,printfolios=true}


\begin{document}

\title[Privacy by Design, Privacy by Expectation]{Privacy by Design, Privacy by Expectation: Aligning Developer Practices and User Expectations in Canadian mHealth Applications}

\author{Opeyemi Ajegbomogun}
\affiliation{%
  \institution{University of Calgary}
  \city{Calgary}
  \country{Canada}}
\email{opeyemi.ajegbomogun@ucalgary.ca}

\renewcommand{\shortauthors}{Ajegbomogun}

\begin{abstract}
Mobile health (mHealth) applications now mediate care for millions of people,
collecting uniquely sensitive information such as medication schedules, mood
logs, and location data, often under consumer-app privacy norms rather than
clinical-grade safeguards. Prior work has repeatedly found troubling gaps
between what these apps claim and what they actually do: flawed consent and
deletion mechanisms, covert data collection in women's health
apps~\cite{malki2024exploring}, and third-party SDK risks that persist even
as platform policies improve, developers often relying on app-market
rejection, not understanding their SDKs'
behavior~\cite{alomar2025effect,alomar2022developers}. Most of this evidence
comes from one vantage point at a time: auditing technical behavior and
policy text, or surveying only one side of the relationship, developers or
users, rarely both, and almost always within the U.S.\ or EU regulatory
context privacy research still defaults to~\cite{sharma2026local}. This
leaves a practice-relevant question: when
developers believe they are building a privacy-respecting app, does that
match what the people actually using it expect, and does either account
match what the app is observed to do? The closest analogue, a recent study
of military-marketed apps, pairs technical analysis with a user study and
finds clear misalignment, but without bringing developers' self-reported
intent into the comparison~\cite{shinkle2026noprivacy}. This study addresses
that remaining gap for the Canadian mHealth ecosystem, governed by PIPEDA
rather than GDPR, through a triangulated, mixed-methods design spanning four
sources: a structured developer survey on reported privacy-by-design
practices and the factors that help or hinder adopting them; a structured
user survey and follow-up focus group on Canadian users' privacy
expectations, concerns, and preferences; a thematic review of selected
apps' privacy policies; and a corpus of roughly 1.45 million app store
reviews, filtered and classified at scale using a natural-language-inference-and-large-language-model
pipeline adapted from recent software-engineering
research~\cite{ansari2026coseto}, surfacing privacy- and security-relevant
experiences, defined against a reconciled privacy
taxonomy~\cite{solove2006taxonomy,deng2011linddun}, that surveys alone would
miss. Developer and user data are coded against a jointly developed
thematic codebook, enabling direct comparison across sources rather than
separate, siloed analyses. We anticipate both convergence, such as shared
concern over opaque third-party data sharing, and meaningful divergence,
particularly where developers consider informal or architectural privacy
measures sufficient but users do not perceive or trust them as such. The
contribution is twofold: an empirically grounded account of where developer
intent and user expectation align or diverge in an under-studied regulatory
context, and a reusable methodology for triangulating self-reported and
observed privacy evidence across stakeholder groups in mobile health
research.
\end{abstract}

\keywords{mHealth, privacy by design, user privacy expectations, Canadian privacy law, PIPEDA, app store reviews, thematic analysis}

\maketitle

\section{Target Journals}
\label{sec:journals}

Table~\ref{tab:journals} compares three venues at the intersection of
software engineering, privacy, and security that fit this work's scope.

\begin{table}[h]
  \caption{Journal Selection Matrix}
  \label{tab:journals}
  \small
  \begin{tabular}{p{0.16\linewidth}p{0.24\linewidth}p{0.24\linewidth}p{0.24\linewidth}}
    \toprule
     & \textbf{IEEE TDSC} & \textbf{ACM TOPS} & \textbf{PoPETs (Primary)} \\
    \midrule
    Full name &
      IEEE Trans.\ on Dependable and Secure Computing &
      ACM Trans.\ on Privacy and Security &
      Proceedings on Privacy Enhancing Technologies \\
    \addlinespace
    Scope &
      Dependability, fault tolerance, and security of computing systems &
      Privacy and security research spanning systems, crypto, and HCI &
      Practical and theoretical privacy-enhancing technologies, real-world systems, privacy engineering \\
    \addlinespace
    Access model &
      Hybrid; optional OA fee &
      Hybrid; optional OA fee &
      Diamond open access (free to publish, free to read) \\
    \addlinespace
    Review cycle &
      Continuous submission; several months &
      Continuous submission; several months, revise-and-resubmit &
      Quarterly; decision within $\sim$2 months; Accept/Revise/Reject \\
    \bottomrule
  \end{tabular}
\end{table}

\subsection{Primary Target: Proceedings on Privacy Enhancing Technologies (PoPETs)}

PoPETs publishes novel, practical, and theoretical research into
privacy-enhancing technologies, real-world systems applications, and
software engineering privacy frameworks, making it a strong fit for a study
that triangulates developer practice, user expectation, and observed
behavior. Unlike IEEE TDSC and ACM TOPS, which remain hybrid subscription
venues with optional author-pays open access, PoPETs is Diamond Open
Access, fully funded by the PETS community, removing both the paywall and
the open-access fee barrier. Its quarterly submission model offers a
predictable, fast review cycle, with decisions typically within two months
and standard journal-style revision options (Accept, Revise, or Reject),
which fits a graduation or project timeline better than a continuous,
open-ended review process.

Per the official PoPETs author guidelines, the abstract must appear
unindexed on the first page as a single, cohesive paragraph within the
\texttt{abstract} environment of the mandatory PoPETs \LaTeX\ template
(satisfied above), written for a broad, non-specialist computer science
audience, and the first page must explicitly connect the work to
real-world privacy applications, which this abstract does by grounding the
research gap in concrete, cited findings from the mHealth privacy
literature rather than an abstract framing of the problem.

\subsection{Citation Style}

PoPETs uses a numeric citation style within an ACM/IEEE hybrid
presentation format, implemented here via
\texttt{ACM-Reference-Format}\slash \texttt{acmart}, as already configured
in this template. Submissions undergo double-blind review, so self-cited
prior work must be referenced in the third person rather than as ``our
previous work.'' There is no page limit on the bibliography; comprehensive
citation is expected.

%%
%% The next two lines define the bibliography style to be used, and
%% the bibliography file.
\bibliographystyle{ACM-Reference-Format}
\bibliography{references}

\end{document}
